\documentclass[11pt]{article}
\usepackage{amsmath,amssymb}
\title{An exact derivative-network cancellation and its mixed-boundary defect}\author{}\date{}
\begin{document}\maketitle
\paragraph{Input and the derivative issue.}
The common-gauge current estimate of the mixed grid joint current with connectors calculation holds for its stated bounded
partial-permutation connectors. In the reflected-current case the
endpoint product in its integration-by-parts formula is indispensable;
its independent coordinate moment estimate pays that product. Piecewise
substochasticity and adjoint connectors justify its mixed loop bound.
An arbitrary nonlocal projection or an unverified actual vertex connector
is still outside that lemma.

The desired additional factor $1/(1+T)$ must not be inferred merely by
counting two derivatives. On unequal proper times, derivatives on a short
propagator generate short-time contact terms. This note proves their
EXACT cancellation for a fully contracted scalar derivative theta, and
locates the surviving defect when its neutral component grid changes.
This is a reusable original-operator identity, not a bound on the entire
Wilson vector/ghost coefficient.

\paragraph{1. Keep the original killing boundary term.}
On one original mixed scalar grid write
\[
 (L_\epsilon u)(x)=\sum_{y\sim x}c_{xy}
 [u(x)-e^{i\epsilon a_{xy}}u(y)]+k_xu(x),\qquad c_{xy}=c_{yx}.
\]
Here $k_x$ is the original Dirichlet killing rate. Suppressed Neumann
outward edges contribute neither a hopping nor a killing term.
For oppositely charged functions define
\[
 \Gamma_\epsilon(u,v)(x)=\sum_{y\sim x}c_{xy}
 [u(x)-e^{i\epsilon a_{xy}}u(y)]
 [v(x)-e^{-i\epsilon a_{xy}}v(y)]+k_xu(x)v(x).                    \tag{1}
\]
Direct multiplication, including the ONE killing term in (1), gives
\[
 \Gamma_\epsilon(u,v)=vL_\epsilon u+uL_{-\epsilon}v-L_0(uv).
                                                                    \tag{2}
\]
Thus for any neutral weight $w$ and the original counting measure,
\[
 \sum_x w\Gamma_\epsilon(u,v)=
 \sum_x[w v L_\epsilon u+w u L_{-\epsilon}v-uvL_0w].              \tag{3}
\]
There is no discarded boundary summation term: the killing term accounts
for it exactly. Equation(3) is a contracted sum over ALL scalar gradient
directions; a particular Wilson curl tensor is not silently substituted
for it.

\paragraph{2. The complete same-grid quadratic flux response vanishes.}
Let $G_\epsilon=L_\epsilon^{-1}$ and $G_0=L_0^{-1}$. Define the actual
contracted scalar network
\[
 \mathcal D(\epsilon)=\sum_{x,z}G_0(x,z)
 \Gamma_\epsilon(G_\epsilon(\cdot,z),G_{-\epsilon}(\cdot,z))(x).
\]
Using $L_\epsilon G_\epsilon=I$ in (3) gives the exact resolvent identity
\[
 \boxed{\mathcal D(\epsilon)=\sum_x
 [2G_0(x,x)G_\epsilon(x,x)-G_\epsilon(x,x)^2]
 =\mathcal D(0)-\sum_x[G_\epsilon(x,x)-G_0(x,x)]^2.}              \tag{4}
\]
A Hermitian magnetic inverse has real diagonal, and
$G_{-\epsilon}=\overline{G_\epsilon}$. Its diagonal is therefore even
in $\epsilon$. Consequently
\[
 \boxed{\mathcal D''(0)=0}                                        \tag{5}
\]
for EVERY finite original grid, with no size estimate or asymptotic fit.

Equivalently, for the scalar heat theta $\Theta(t_1,t_2,t_3;\epsilon)$,
the contracted-gradient network equals
$-\partial_{t_1}\Theta-\partial_{t_2}\Theta+\partial_{t_3}\Theta$.
Integrating all three proper times produces the THREE zero-time contacts
in (4). Their flux Hessians cancel. This explains why taking absolute
values of the separate short-time contacts cannot supply the desired
marginal derivative estimate: the signed contacts must first be combined.
No numerical $1/T$ factor was inserted into a propagator.

\paragraph{3. A different original neutral grid produces a boundary defect.}
Let $K_0$ be another original scalar component or ghost operator, and let
$P$ be a fixed bounded translation/overlap partial permutation mapping its
grid into the charged grid. Keep its true inverse and put
\[
 W=P K_0^{-1}P^*,\qquad E=L_0 W-I.
\]
Repeat (3) with the neutral column $w=W(\cdot,z)$. The exact answer is
\[
 \boxed{\mathcal D_W(\epsilon)=
 2\sum_x W(x,x)G_\epsilon(x,x)-\sum_xG_\epsilon(x,x)^2
             -\sum_{x,z}E(x,z)|G_\epsilon(x,z)|^2.}               \tag{6}
\]
In particular, with dots denoting background derivatives at zero,
\begin{align}
 \mathcal D_W''(0)
 &=2\sum_x[W(x,x)-G_0(x,x)]\ddot G(x,x)
             -\sum_{x,z}E(x,z)(|G_\epsilon(x,z)|^2)''_0\\
 &=\sum_{u,z}E(u,z)
 [2G_0(z,u)\ddot G(z,z)-(|G_\epsilon(u,z)|^2)''_0].              \tag{7}
\end{align}
The second line uses $W-G_0=G_0E$. The diagonal $u=z$ summand vanishes.
The whole bracket is gauge invariant; expanding only one of its current
terms would lose that property.

The defect has the literal operator decomposition
\[
 E=(PP^*-I)+(L_0P-PK_0)K_0^{-1}P^*.                              \tag{8}
\]
For the original component grids and a bounded coordinate translation,
the two flat difference operators agree in their common interior.
Consequently $L_0P-PK_0$ has rows only in a fixed-width boundary collar;
$PP^*-I$ is supported in the missing overlap collar. Thus the first index
$u$ in (7) lies in $O(R^3)$ boundary rows, rather than an unrestricted
$O(R^4)$ volume set. Its second index remains unrestricted and its actual
kernel $K_0^{-1}$ is still present. This is a SUPPORT cancellation, not
an estimate of the defect's weighted spatial moment.

If a proposed vertex uses a parameter-dependent connector, (6) still
holds pointwise with that connector, but its derivatives also act on
$W$ and $E$. Those partners are NOT omitted or covered by (7)'s fixed-$P$
formula. Long tree connectors and nonlocal projectors are likewise not
replaced by the bounded overlap map used here.

\paragraph{4. Direct original-grid checks.}
On the $R=3$ grid with reflected direction one, direct complex inversion
checks (4) through $8\,10^{-16}$. Its deficit scales as $\epsilon^4$;
$[\mathcal D(0)-\mathcal D(\epsilon)]/\epsilon^4$ approaches
$1.356\,10^{-9}$ in that particular finite test.
An independent $R=4$ check uses direction-one and direction-zero grids
of 108 sites each, with their literal overlap map. The mixed identity(6)
passes through $7\,10^{-15}$, and the defect on strict common-interior rows
is below $10^{-15}$. These tests include the original killing diagonals.

\paragraph{Precisely what remains unpaid.}
Equations(4)--(8) cancel the entire bulk scalar contraction and retain its
mixed-boundary remainder. To use them for the complete first Wilson
coefficient, the actual curl/ghost tensor contractions and differentiated
endpoint/spin/projector terms must first be grouped into this scalar
energy contraction plus explicitly retained remainders. That grouping has
not been established here. Even after it, the weighted boundary defect(7)
needs a physical-energy estimate consistent with the finite boundary nonlinear ward cancellation calculation's complete
boundary Ward. No bound on the full $\mathcal C_R''$ or on the original
growing-block differentiated remainder follows merely from (5).
No S94 or physical-scale induction is claimed.
\end{document}
